\section{A strict crossed Laurent box with a full-original-tester compact capacity gap}
\label{sec:r16-crossed-box-compact-gap}

\paragraph{Original category and the statement being tested.}
The construction uses the prescribed finite extension
$K/\mathbf Q_p$, its unchanged absolute value, and every original
classical normalized tester. It does not replace the tester class by
toric functions. A toric majorant will be used only to obtain an
\emph{upper bound at a specified Gauss atom}, by the frozen original
contact inequality. In particular no nonlinear pushforward formula
for arbitrary psh functions is assumed.

Fix a uniformizer $\pi$, and put
\[
       r=|\pi|,\qquad \lambda=-\log r>0,\qquad
       R=r^{-1},\qquad \kappa=(2\lambda)^{-q},
       \qquad q\ge2.
 \tag{CB1}
\]
Let
$X=(\mathbf A^{2q}_{K,X_1,\ldots,X_q,Y_1,\ldots,Y_q})^{\rm an}$.
Inside the compact radius-$R^3$ original affinoid polydisc, use
the strict Laurent box
\[
 \begin{split}
 V=\{\, &|X_i|<R,\ |Y_j|<R,\ |X_iY_j|<R
                         \quad(1\le i,j\le q)\,\},\\
 K_0=\{\, &|X_i|\le R,\ |Y_j|\le R,\ |X_iY_j|\le R
                         \quad(1\le i,j\le q)\,\}.
 \end{split}
 \tag{CB2}
\]
Thus $K_0\Subset\operatorname{Int}_X(B)$ for that affinoid $B$.
All the functions defining the upper inequalities are original
analytic functions and all thresholds are in $|K^\times|$. No
lower inequalities are needed for this allowed strict Laurent box.
The space $X$ is good, Hausdorff, boundaryless, quasi-smooth and
pure of dimension $2q$.

Write $D=d'd''$, and define the original full cap currents
\[
 \begin{split}
 S_1&=\bigwedge_{j=1}^q
                   D\max\{\log|Y_j|,-2\lambda\}\wedge[X],\\
 S_2&=\bigwedge_{i=1}^q
                   D\max\{\log|X_i|,0\}\wedge[X],\\
 S&=2S_1+S_2.
 \end{split}
 \tag{CB3}
\]
Every logarithmic hinge is locally constant across its zero
divisor. These are original locally bounded positive cap products,
well defined on an ordinary neighborhood of $K_0$. They are
symmetric middle-positive of bidegree $(q,q)$ and separately
$d'$- and $d''$-closed. Their stronger special construction is not
a change in the positivity hypothesis of the statement refuted.

The R9 full-cap support and Fubini theorem identifies the carriers as
\[
 F_1=\operatorname{pr}_Y^{-1}(\zeta_{r^2,\ldots,r^2})\cap V,
 \qquad
 F_2=\operatorname{pr}_X^{-1}(\zeta_{1,\ldots,1})\cap V.
 \tag{CB4}
\]
These are the \emph{full fibers of the joint Gauss seminorms}.
They are not the entire coordinate norm levels, nor an intersection
of coordinate conditions with correlated points left unexamined.
On $F_1$ the free $X$ domain is the full radius-$R$ open
polydisc. On $F_2$ the free $Y$ domain is the same full open
polydisc. Indeed on the former carrier $|Y_j|=r^2$, and on the
latter $|X_i|=1$, so the product inequalities in (CB2) introduce
no additional free-coordinate restriction. The original cap
calibration has normal multiplicity one.

For every original classical psh $a:V\to[0,1]$, put
\[
 \mu_{i,a}=(Da)^q\wedge S_i,\qquad
 \mu_a=(Da)^q\wedge S=2\mu_{1,a}+\mu_{2,a}.
 \tag{CB5}
\]
All these are the original positive Radon top measures, on the
complete original test sheaf and every actual refinement. Define
\[
 C(E)=\sup_{a:V\to[0,1]\ {\rm original\ classical\ psh}}\mu_a(E),
 \qquad
 C^*(E)=\inf_{\substack{G\supset E\\G\ {\rm ordinary\ open\ in}\ V}}C(G).
 \tag{CB6}
\]
No boundary value is imposed on a normalized capacity tester.
No member of this family is silently placed in a zero-boundary
energy completion.

\subsection{The three original Gauss points}

Let $A,B,z_0$ denote the original joint product Gauss points with
the following radii:
\[
 \begin{array}{c|c|c}
       &X_1,\ldots,X_q&Y_1,\ldots,Y_q\\ \hline
 A     &(r,\ldots,r)&(r^2,\ldots,r^2)\\
 B     &(1,\ldots,1)&(1,r,\ldots,r)\\
 z_0   &(1,\ldots,1)&(r^2,\ldots,r^2).
 \end{array}
 \tag{CB7}
\]
All lie in the compact unit polydisc inside $V$. Set
$K=\{A,B\}$. The point $A$ belongs to $F_1\setminus F_2$,
and $B$ belongs to $F_2\setminus F_1$; $z_0$ belongs to their
intersection. In particular, the following calculation does not
replace $F_1\cap F_2$ by the singleton $\{z_0\}$.

\begin{theorem}[A compact raw/outer gap for every original tester]
\label{thm:r16-crossed-box-compact-capacity-gap}
For the fixed original field and every $q\ge2$, the strict box
and allowed background in (CB1)--(CB3) satisfy
\[
 \boxed{\qquad
 C(K)=\bigl(2+2^{1-q}\bigr)(2\lambda)^{-q}
       \ <\ 3(2\lambda)^{-q}\ \le C^*(K)<\infty .
 \qquad}
 \tag{CB8}
\]
The raw supremum on $K$ is attained by an explicit actual
continuous original normalized tester. There is also a sequence
of such actual normalized testers whose source measures converge
weakly on the original space to a measure with mass
$3(2\lambda)^{-q}$ on $K$.

Consequently the full general compact no-gain assertion FG1 is
\texttt{FALSE}, already on a strict Laurent box and a separately
closed allowed background extending across its closed box. The
same example refutes universal Borel raw/outer equality. It does
not assert failure for every background or every strict box.
\end{theorem}

\noindent\textbf{Status: \texttt{NEW\_PROVED}} for the
explicit full-original-tester compact gap, subject to the exact
source-hash audits in the companion ledger. The polynomial moving
atom calculation used below is the frozen R3 input, not a newly
claimed example. The new ingredients are the crossed strict box,
the asymmetric full cap pair, and the upper bound for every
actual original tester.

\subsection{Gauss majorization and an atomic comparison}

We prove the full-test comparison needed for (CB8), including
the use of auxiliary fields. They appear only in the
\emph{defining scalar curve tests} of an original classical
function; no Bedford--Taylor theory over an auxiliary field is
introduced.

\begin{lemma}[A full-cap Gauss majorant controls the Gauss atom]
\label{lem:r16-crossed-box-Gauss-atom}
Let $a:V\to[0,1]$ be any original classical psh function, possibly
discontinuous. Fix either carrier in (CB4), write $Z$ for its
$q$ free coordinates, and let $b$ denote the fixed normal joint
Gauss point. For $s\in(-\infty,\lambda)^q$, write $\sigma_b(s)$
for the original joint Gauss point with free radii $e^{s_i}$
and normal point $b$, and put
\[
                         \varphi(s)=a(\sigma_b(s)).
 \tag{CB9}
\]
Then $\varphi$ is finite, continuous, convex and coordinatewise
nondecreasing, with $0\le\varphi\le1$. On the entire full carrier,
at points where all free coordinates are nonzero,
\[
                  a(x)\le
                \varphi(\log|Z_1(x)|,\ldots,\log|Z_q(x)|).
 \tag{CB10}
\]
For every $s$ in the displayed real domain, one has
\[
                   \mu_{i,a}\{\sigma_b(s)\}
                       \le q!\operatorname{Vol}_q\partial\varphi(s).
 \tag{CB11}
\]
Lebesgue volume is in the original logarithmic coordinates and
their standard integral lattice. In particular there is no
unrecorded absolute-value or residue-degree rescaling.
\end{lemma}

\begin{proof}
First prove (CB10). For an arbitrary point of the full carrier,
pass to a complete extension containing its completed residue
field, only to apply scalar curve tests. Hold all but one free
coordinate fixed and replace that coordinate by the variable on
the closed scalar disc of radius equal to its current norm.
The entire closed disc maps into $V$: its radius is strictly
less than $R$, and the fixed normal radii make all product
inequalities automatic. The pulled-back original function is
bounded subharmonic. Its value is at most its value at the
unique boundary Gauss point, by the scalar subharmonic maximum
principle on a closed disc with a slightly larger open disc
available. This principle applies to upper-semicontinuous
subharmonic functions, without a continuity hypothesis.

Repeat for all free coordinates, using the corresponding
completed residue fields if necessary. At the end every free
coordinate has been made Gauss-independent of the unchanged
normal base. For an original polynomial
$\sum_I c_I(Y)Z^I$, its resulting seminorm is exactly
$\max_I|c_I(b)|e^{I\cdot s}$. Thus its original image is
$\sigma_b(s)$. Uniform polynomial approximation on the compact
polydiscs gives the same identity for analytic functions.
This proves majorization on the \emph{entire} carrier, including
points with correlations before the free coordinates are replaced.
It is not a statement just on a chosen skeleton.

We next prove joint convexity. Fix an integral direction
$m\in\mathbf Z^q$ and a segment $s^0+t m$ in the real domain.
Over a complete extension containing the normal Gauss field,
choose independent Gauss coefficients $c_i$ of radii $e^{s_i^0}$.
Use the scalar annulus map
\[
                         Z_i=c_iT^{m_i}.
 \tag{CB12}
\]
At a scalar Gauss point of radius $e^t$, its original image
is precisely $\sigma_b(s^0+t m)$. Indeed, after grouping a
polynomial by the exponent $I\cdot m$ of $T$, the independent
Gauss coefficient norm in each group remains the maximum of
the individual monomial norms. The resulting norm is therefore
$\max_I |c_I(b)|e^{I\cdot(s^0+t m)}$. Negative exponents of
$T$ cause no problem on this scalar annulus. The original image
of the entire annulus stays in $V$ along the chosen segment.

Subtracting the affine harmonic interpolation of two boundary
values and applying the scalar annular maximum principle gives
convexity along that segment. The same conclusion holds for every
rational direction after rescaling. Separate convexity and the
bound $0\le\varphi\le1$ give a local Lipschitz bound on every
smaller real box: each coordinate secant slope is bounded by the
reciprocal distance to the boundary of a slightly larger box.
Changing one coordinate at a time proves interior continuity.
Approximating a real direction by rational directions now proves
full convexity. Finally a convex function of one coordinate on
$(-\infty,\lambda)$ which is bounded above cannot have a
negative secant slope. Each coordinate section is therefore
nondecreasing.

On an ordinary annular neighborhood of $\sigma_b(s)$ define
the original toric function
\[
                           h=\varphi(\log|Z_1|,\ldots,\log|Z_q|).
 \tag{CB13}
\]
Finite convex functions of original unit moments are continuous
classical psh by the frozen original convex-composition
calculus. Only this local annular neighborhood is needed; no
extension across zero coordinates is required. On its intersection
with the full carrier, (CB10) gives $a\le h$, and at
$\sigma_b(s)$ there is equality.

Here is the carrier-order bridge for the contact comparison.
Let $w=\max(a,h)$ on that neighborhood. The functions $w$ and
$h$ agree pointwise on the whole support of $S_i$. Their first
products agree because $(w-h)S_i=0$. Inductively every bounded
BT product is still supported there, and multiplication by
$w-h$ vanishes on its coefficient measures. The defining
recursion therefore gives
\[
                             \mu_{i,w}=\mu_{i,h}.
 \tag{CB14}
\]
This is equality on every original compact test and actual
refinement, using the fixed original background. It is not a
new restriction or division theorem for arbitrary currents.

The frozen R3 maximum inequality, with its strict Borel locality
proof, says
$\mu_{i,\max(a,h)}\ge
\mathbf1_{\{a\ge h\}}\mu_{i,a}$.
At the displayed Gauss point $a=h$, so (CB14) yields
\[
                     \mu_{i,a}\{\sigma_b(s)\}
                            \le\mu_{i,h}\{\sigma_b(s)\}.
 \tag{CB15}
\]
Ordinary order $a\le h$ off the carrier was neither needed nor
assumed.

For the \emph{toric function $h$}, original full-cap calibration
and Fubini give
$\mu_{i,h}\{\sigma_b(s)\}
=q!\operatorname{Vol}_q\partial\varphi(s)$.
The formula holds first for finite convex PL functions: the
top-degree corner mass is $q!$ times the slope-polytope volume,
and the normal cap has multiplicity one. Approximate a general
finite convex $\varphi$ uniformly on a protected real annular
stage by convex PL functions. The frozen bounded uniform BT
continuity, the original toric calibration on every compact
test, and Alexandrov weak continuity pass the full measure
identity to the limit. Equivalently the real measure is
$q!\operatorname{MA}_{\mathbf R}(\varphi)$ carried by the
continuous original joint Gauss section. Its atomic mass is the
displayed subgradient volume. This justifies (CB11) without any
pushforward equality for the original, possibly non-toric
function $a$.
\end{proof}

\paragraph{Exact antecedents for the lemma.}
The scalar disc and annulus maximum principles are the same
classical curve facts used in frozen R14 NT2--NT3 and in R16
WQ7. Joint Gauss convexity for every bounded original tester is
already frozen R6,
\texttt{lem:r6-joint-gauss-convexity}, in
\texttt{capacity\_gauss\_convex\_reduction.tex}. The proof above
spells out its scalar-curve use on the present full caps; joint
convexity by itself is not claimed anew. The frozen R4 definition
of the actual classical class explicitly permits every analytic
curve after every complete isometric extension. Endpoint radii
may be realized in a further such extension when using closed
scalar annuli. The Gauss seminorm/Fubini identities are frozen R9.
The maximum inequality is frozen
\texttt{prop:r3-relative-contact-inequality}, included verbatim
in the R4 aggregate at the source
\texttt{base\_r3/sections/relative\_contact\_inequality.tex}.
Its proof uses strict Borel maximum locality and compact
Portmanteau after a decreasing maximum, not FG1. The toric
point mass normalization is the frozen R3 moving-atom
calculation. Only that toric calibration, and not its already
disproved extension to arbitrary nonlinear Gauss pushforwards,
is used in (CB15).

For primary scalar references see V.~Wanner,
\emph{Subharmonic functions and real-valued differential forms on
non-archimedean curves} (2019), Definition~3.1.5 and
Propositions~3.1.8--3.1.10,
\url{https://epub.uni-regensburg.de/40530/1/THESIS.pdf}.
The auxiliary fields for these scalar tests may be chosen complete
and algebraically closed; no current is transported to them.
For real Alexandrov measure and its weak continuity see
G.~De Philippis and A.~Figalli,
\emph{The Monge--Ampere equation and its link to optimal
transportation}, Section~2.1, formula~(2.2) and Proposition~2.4,
\url{https://people.math.ethz.ch/~afigalli/lecture-notes-pdf/The-Monge-Ampere-equation-and-its-link-to-optimal-transportation.pdf}.
Both primary statements were read directly for this audit.

\subsection{The upper bound for every original normalized tester}

Fix an arbitrary $a$ in (CB6), and put
\[
           t=a(z_0)\in[0,1],\qquad
           m_A=\mu_{1,a}\{A\},\qquad m_B=\mu_{2,a}\{B\}.
 \tag{CB16}
\]
Apply the lemma to $F_1$. The free-coordinate trace
$\varphi_X$ is defined on $(-\infty,\lambda)^q$; put
$v=\varphi_X(-\lambda,\ldots,-\lambda)\ge0$.
For every $p\in\partial\varphi_X(-\lambda,\ldots,-\lambda)$,
coordinate monotonicity gives $p_i\ge0$. The supporting-plane
inequality, first at the point $(0,\ldots,0)$ and then at
$(\lambda-\epsilon,\ldots,\lambda-\epsilon)$, gives
\[
        \lambda\sum_i p_i\le t-v,\qquad
        2\lambda\sum_i p_i\le1-v
                        \quad\hbox{after }\epsilon\downarrow0.
 \tag{CB17}
\]
The point with all free radii one and this normal base is
\emph{exactly} $z_0$. It is not the retraction of a different
residue-branch point. This exact choice of the asymmetric cap
radii is essential.

The nonnegative simplex
$\{p:\sum_i p_i\le c\}$ has volume $c^q/q!$.
Using (CB11), (CB17), and $v\ge0$ therefore yields
\[
                  m_A\le
                 \lambda^{-q}\min\{2^{-q},t^q\}.
 \tag{CB18}
\]

For $F_2$, let $\varphi_Y$ be the free $Y$ trace, and put
$b=(0,-\lambda,\ldots,-\lambda)$.
Coordinate monotonicity gives
\[
       \varphi_Y(b)\ge
          \varphi_Y(-2\lambda,\ldots,-2\lambda)=t.
 \tag{CB19}
\]
For $p\in\partial\varphi_Y(b)$ one has $p_i\ge0$, and
approaching the all-upper corner gives
\[
       \lambda p_1+2\lambda\sum_{i=2}^q p_i
                           \le1-\varphi_Y(b)\le1-t.
 \tag{CB20}
\]
This simplex has volume
$(1-t)^q/(q!\,2^{q-1}\lambda^q)$, so (CB11) gives
\[
                    m_B\le
                    \frac{(1-t)^q}{2^{q-1}\lambda^q}.
 \tag{CB21}
\]

Since the two target points are off the opposite carrier,
$\mu_a(K)=2m_A+m_B$. Multiplying the last bounds by
$(2\lambda)^q$ gives
\[
       (2\lambda)^q\mu_a(K)
          \le2\min\{1,(2t)^q\}+2(1-t)^q.
 \tag{CB22}
\]
On $[0,1/2]$ the right side is a convex function of $t$,
so its maximum is at an endpoint. Its values there are
$2$ and $2+2^{1-q}$. On $[1/2,1]$ it equals
$2+2(1-t)^q$ and is nonincreasing. Hence, for
\emph{every} original normalized tester,
\[
                       \mu_a(K)\le
                           (2+2^{1-q})(2\lambda)^{-q}.
 \tag{CB23}
\]
This bound includes all non-toric and all discontinuous
classical testers. No restricted tester supremum occurs in it.

\subsection{An actual tester attaining the raw upper bound}

Put $L_X=\max_i\log|X_i|$, with the value $-\infty$ if every
$X_i$ is zero. Define
\[
 a_*=\max\left\{
 0,\ \frac{L_X+\lambda}{2\lambda},\
 \frac{L_X+\log|Y_1|+\lambda}{2\lambda},\
 \frac{L_X+\log|Y_l|+3\lambda}{4\lambda}\ (2\le l\le q)
                \right\}.
 \tag{CB24}
\]
After distributing the finite maxima, every branch is a positive
multiple of the logarithm of an original coordinate or coordinate
product, plus a constant. The finite floor removes the poles.
Thus $a_*$ is a continuous original classical psh function.
Every nonconstant branch is strictly below one on $V$:
the first uses $|X_i|<R$ and the others use $|X_iY_l|<R$.
Consequently $0\le a_*\le1$ on the whole original domain.

On the full $F_1$, where $\log|Y_l|=-2\lambda$, its trace is
\[
                    a_*|_{F_1}
                        =\max\{0,(L_X+\lambda)/(2\lambda)\}.
 \tag{CB25}
\]
For the branches with $l\ge2$ the expression is one half of
$(L_X+\lambda)/(2\lambda)$, which is dominated by that
expression or by zero. The $Y_1$ branch is smaller still.
Original full-cap calibration therefore gives
$\mu_{1,a_*}=\kappa\delta_A$.

On the full $F_2$ one has $L_X=0$, so
\[
 a_*|_{F_2}=
 \max\left\{\frac12,\
       \frac{\log|Y_1|+\lambda}{2\lambda},\
       \frac{\log|Y_l|+3\lambda}{4\lambda}\ (2\le l\le q)
       \right\}.
 \tag{CB26}
\]
The constant and all $q$ free coordinate branches tie at the
single real logarithmic vertex $(0,-\lambda,\ldots,-\lambda)$.
Its full original cap calibration puts the top-degree mass at
the joint Gauss point $B$, not on that vertex's entire norm level.
The slope simplex gives mass
\[
       \frac1{2\lambda(4\lambda)^{q-1}}
                          =2^{1-q}\kappa.
 \tag{CB27}
\]
On every other real logarithmic face the active slope differences
have rank less than $q$, so there is no other top mass. The full
original cap calculation gives
\[
               \mu_{a_*}
                  =\kappa(2\delta_A+2^{1-q}\delta_B).
 \tag{CB28}
\]
Combining (CB23) and (CB28) proves the exact raw equality in
(CB8).

\subsection{Actual moving sources force a larger outer mass}

Choose pairwise distinct monic irreducible residue polynomials
$\overline P_j$ of degrees $d_j\ge2$, and monic integral lifts
$P_j$. Let $\xi_j$ be the unique original point above the
radius-$r^{d_j}$ Gauss point for $P_j$ in its closed residue
branch. This is exactly the frozen R3 moving-atom construction.
Its original residual normal degree is $d_j$.

Let $B_j$ be the original product point with $X$ radii all
one, first $Y$ factor $\xi_j$, and other $Y$ radii $r$.
Then
\[
                            B_j\longrightarrow B.
 \tag{CB29}
\]
Indeed a fixed polynomial has only finitely many irreducible
residue factors, so eventually its norm at $\xi_j$ is its
unit Gauss norm. The same coefficientwise Gauss calculation
and uniform polynomial approximation apply after adjoining the
other fixed Gauss coordinates. This gives convergence in the
original Berkovich topology. The sequence and its limit lie
in the compact unit polydisc inside $V$.

Use the finite-logarithm convention and put
\[
 \begin{split}
 L_j(Y)&=\max\{d_j^{-1}\log|P_j(Y_1)|,\
                       \log|Y_2|,\ldots,\log|Y_q|\},\\
 U_j&=\max\{L_X+L_j-\lambda,\
              L_X-2\lambda,\ L_j-2\lambda,\ -2\lambda\},\\
 a_j&=1+U_j/(2\lambda).
 \end{split}
 \tag{CB30}
\]
Again these are actual continuous finite maxima of original
logarithmic psh functions with a finite floor. We check their
\emph{global} normalization rather than only their cap traces.
Write $\rho_X=\max_i\log|X_i|$ and
$\rho_Y=\max_l\log|Y_l|$. Monicity and integrality give
\[
                    L_j\le\max\{0,\rho_Y\},\qquad L_X=\rho_X.
 \tag{CB31}
\]
If $\rho_Y\le0$, then $L_X+L_j<\lambda$ by the individual
$X$ inequalities. If $\rho_Y>0$, the cross-product inequalities
give $\rho_X+\rho_Y<\lambda$, again proving
$L_X+L_j<\lambda$. The other branches in (CB30) are also
nonpositive, because $L_X<\lambda$ and $L_j<\lambda$.
The floor gives
\[
                              -2\lambda\le U_j\le0,
                   \qquad 0\le a_j\le1 \quad\hbox{on all }V.
 \tag{CB32}
\]

On $F_1$, the normal $Y_1$ radius is $r^2<1$. The constant
coefficient of $P_j$ is a unit, so $|P_j(Y_1)|=1$ on that
entire fiber. Consequently $L_j=0$ and
\[
                 U_j|_{F_1}
                    =\max\{L_X-\lambda,-2\lambda\}.
 \tag{CB33}
\]
Its unit original top mass is at $A$. On $F_2$, one has
$L_X=0$, and hence
\[
                 U_j|_{F_2}
                     =\max\{L_j-\lambda,-2\lambda\}.
 \tag{CB34}
\]
This is exactly the frozen R3 function $g_j$ in its free
$Y$ variables. The slope $1/d_j$ cancels the original residual
degree $d_j$, so its full original source is the unit
$\delta_{B_j}$. The R9 full-cap trace and Fubini identities
therefore give, on every original test and actual refinement,
\[
                       \mu_{a_j}
                           =\kappa(2\delta_A+\delta_{B_j}).
 \tag{CB35}
\]
These are actual normalized tester measures on the original
crossed domain. No cutoff, completed representative, or
zero-boundary source selection has entered their construction.
The functions also have ordinary-neighborhood psh extensions
across $K_0$, although no energy boundary theorem is needed.

Let $G$ be any ordinary open neighborhood of $K=\{A,B\}$ in
$V$. By (CB29), $B_j\in G$ for every sufficiently large $j$.
Since $A\in G$, (CB35) gives
\[
                             C(G)\ge3\kappa.
 \tag{CB36}
\]
Taking the infimum over all such $G$ yields $C^*(K)\ge3\kappa$.
A relatively compact ordinary neighborhood of $K$ has finite
capacity by local CLN, so $C^*(K)<\infty$. Because $q\ge2$,
$2^{1-q}<1$, and (CB23)--(CB28) prove (CB8), with the explicit
gap
\[
                      C^*(K)-C(K)
                         \ge(1-2^{1-q})(2\lambda)^{-q}>0.
 \tag{CB37}
\]

\subsection{What changes, and what does not follow}

For any $\epsilon$ smaller than the right side of (CB37), no
actual normalized original tester can satisfy
$\mu_a(K)\ge C^*(K)-\epsilon$. Thus this is an explicit
counterexample to the full FG1 quantifiers, rather than merely
nonattainment of a supremum or failure of a preferred
zero-contact normalization. In fact the raw supremum is
attained. Since $K$ is finite and Borel, the universal Borel
raw/outer equality is also false.

The source measures in (CB35) converge weakly to
$\kappa(2\delta_A+\delta_B)$, whose compact mass exceeds the
raw capacity. This realizes the no-gain failure in the exact
localized-weak-cluster language of frozen CE5. CW's literal
strong-data conclusion for localized clusters is not
contradicted. Its conclusion concerns the energy datum class,
not recovery by a normalized actual potential.

The domain geometry is essential. On the rectangular polydisc
used in WC, the product branches in (CB24) and (CB30) do not
have the same global normalization. The strict crossed box
imposes exactly the original product bounds required here,
while each full cap still sees the entire free polydisc.
No original-field extension or change of absolute value is
responsible for the gap.

The obstruction is nonlinear: at $q=1$ the numerical gap in
(CB37) is zero. No claim is made against a surviving
degree-one or other restricted compact regularity theorem.
No positive statement about arbitrary $S$ is obtained by
extrapolating this special background.

The arbitrary actual-analytic-zero KC12/GQ/CTO problem,
intrinsic geometric criteria for completed Dirichlet
compatibility, and literal positive ED-to-strong continuity
are not resolved by this construction. They retain
\texttt{OPEN} status. Every energy theorem used elsewhere
retains its original strict-box, no-loss and boundary
conditions; no capacity tester here is declared a completed
zero-boundary source merely because it is bounded.

\paragraph{Errors and rejected shortcuts.}
The earlier symmetric-cap choice with a point in a different
residue branch did not identify its actual value with the
toric Gauss trace. A toric upper bound at that different point
could not furnish the needed lower bound. The asymmetric
choice in (CB7) corrects this: the $X$-unit Gauss point is
literally the same $z_0$ used by both carriers. The proof
does not use exact saturation rigidity.

It would also be invalid to identify the nonlinear curvature
of an arbitrary tester with the real Hessian measure of its
Gauss trace. Only the toric majorant has that calibration.
The contact inequality and the full-carrier equality (CB14)
provide the required one-sided atomic comparison. Likewise,
the cap carriers cannot be replaced by their whole norm
levels, and the normal residue factor $d_j$ in (CB34) cannot
be discarded.

\paragraph{Provenance and status ledger.}
The capacity researcher constructed the crossed-domain
normalization, the asymmetric pair of full caps and the
two-point tradeoff. The root independently identified the
residue-point/Gauss-retraction mismatch in an earlier
candidate. The GQ and literal researchers independently
reconstructed the scalar-disc/annulus Gauss-majorization
argument and the atomic simplex bounds. The literal
researcher supplied (CB24), and the capacity researcher
checked its exact traces and normalization. These contributions
are disclosed in the source-hash review records.

The original bounded BT, full-cap Fubini, strict Borel
maximum/contact inequality, toric calibration and R3 moving
atom calculation are \texttt{INTERNAL\_INPUT}. Scalar
subharmonic maximum principles, elementary convex
subgradients and Alexandrov weak continuity are
\texttt{EXTERNAL\_KNOWN}. The all-tester upper bound,
matching raw maximizer and strict compact gap form one
\texttt{NEW\_PROVED} construction; its elementary status
consequences are \texttt{DIRECT\_COROLLARY}. Universal FG1
and universal Borel raw/outer equality are \texttt{FALSE}.
Claims extending a restricted surviving theorem beyond its
printed hypotheses remain \texttt{CONDITIONAL}; KC12/GQ/CTO,
general intrinsic geometric compatibility and literal LG
remain \texttt{OPEN}. The incorrect nonlinear pushforward
and the incorrect residue-point identification are recorded
as \texttt{FAILED\_ATTEMPT}, not omitted from the audit.
